\par\medskip
\noindent\textbf{Register pressure and rendering time in the paper}\par
We matched the five variants in Tables~\ref{tab:hydra-summary-0-maximum}
and~\ref{tab:hydra-summary-0-path-traversal} to Table~5 of
\texttt{crossrt\_ieee\_template.tex} (revision \texttt{1d6010f}). Register
counts are from \texttt{test\_035}. Times are the measurements reported
for Coffee, Spaceship, Staircase and Kitchen at 1024$\times$1024 and 256 spp.
For each GPU and workload, we compare wavefront with megakernel within each
traversal mode, then compare the two hardware megakernels. Each cell gives
the change in peak registers from Table~\ref{tab:hydra-summary-0-maximum}
and the geometric mean of the four scene-wise time ratios
$(\prod_s t_{\mathrm{after},s}/t_{\mathrm{before},s})^{1/4}$.
A ratio below one means shorter render time.

\begin{center}
\begin{tabular}{llccc}
\hline
GPU & Workload & SW wave $\to$ mega & HW wave $\to$ mega & HW mega $\to$ mega-p \\
    &          & $\Delta$reg. / time ratio & $\Delta$reg. / time ratio & $\Delta$reg. / time ratio \\
\hline
RTX 4090   & PT   & $+16$ / 0.806 & $+16$ / 0.430 & $0$ / 0.646 \\
RX 6800 XT & PT   & $+32$ / 1.304 & $+32$ / 1.148 & $+8$ / 0.944 \\
Arc A380   & PT   & $+123$ / 1.071 & $+153$ / 2.450 & $-5$ / 0.929 \\
\hline
RTX 4090   & RTAO & $0$ / 0.998 & $0$ / 0.956 & $0$ / 1.068 \\
RX 6800 XT & RTAO & $+8$ / 0.999 & $+16$ / 1.001 & $0$ / 1.075 \\
Arc A380   & RTAO & $-41$ / 1.200 & $+62$ / 0.771 & $0$ / 4.220 \\
\hline
\end{tabular}
\end{center}

\noindent\textbf{Path tracing.}
Fusion raised the register counts in both tables on all three GPUs in both
traversal modes. With hardware traversal, the megakernel shortened render
time by 57\% on RTX 4090 and increased it by 15\% on RX 6800 XT and 145\% on
Arc A380. The effect had the same sign in all four scenes on NVIDIA and Intel,
and in three on AMD. RADV's reported limit fell from 8 subgroups per SIMD for
the highest-VGPR wavefront kernel (\texttt{kernel\_NextBounce}) to 6 for the
hardware megakernel. On Arc A380, ANV reported 352 spill operations for the
hardware megakernel, versus zero for the standalone hardware traversal stage.
Enabling persistent rays in the hardware PT megakernel shortened the geometric
mean time on all three GPUs while the measured count stayed at 80 on RTX 4090,
rose from 160 to 168 on RX 6800 XT, and fell from 317 to 312 on Arc A380.

\par\medskip
\noindent\textbf{Ambient occlusion.}
RTX 4090 reported a peak of 48 registers for every variant. On RX 6800 XT,
the extra 8--16 VGPRs in the megakernel accompanied less than 0.2\% change in
geometric mean time in either traversal mode. On Arc A380, the software
megakernel used fewer peak registers and took 20\% longer, while the
hardware megakernel used more and took 23\% less time. Both directions held in all
four scenes. The persistent hardware megakernel on Arc A380 kept the count at
116 and took 4.22 times as long. The local persistent capture is pending
render validation as described at the start of this supplement.

\par\medskip
\noindent\textbf{Conclusion.}
In both software and hardware PT, switching from wavefront to megakernel
increased peak register counts on all three GPUs. On average, rendering
slowed down on RX 6800 XT and Arc A380 and sped up on RTX 4090. RTAO and
persistent-ray comparisons showed no consistent relationship. Register counts
and timings come from different scenes, and the variants change several
aspects of execution, so the results establish associations between
configurations.
\par\medskip
